\section{补充结果与审计协议}

\subsection{历史 FB 候选评分的来源边界}

\begin{table}[h]
\caption{FB15K237-2 历史十选一评估（百分比，10,000 条测试问题）。InfoNCE-i0 使用历史候选文件和不同训练配置，不是 SC-CL 的跨数据集结果。}
\label{tab:historical-fb}
\centering\scriptsize
\begin{tabular}{lrrrr}
\toprule
方法 & 官方 EM/H@1 & MRR & H@3 & H@10 \\
\midrule
\bone & 79.91 & 88.01 & 95.80 & 100.00 \\
历史 InfoNCE-i0 & 91.62 & 95.41 & 99.20 & 100.00 \\
\bottomrule
\end{tabular}
\end{table}

两组评分均来自相同的本地 processed test，哈希为
\texttt{c77dde51aafc7b1ce3a5011ade914343c2689f8b38b12774ea4840ec1c9068ed}。
历史 i0 使用 \texttt{fb15k237\_i0\_paired\_candidates.json}，LoRA rank 为 24、Stage 2
训练 5 个 epoch、候选损失权重为 0.1；这些条件与本文 SC-CL 的训练契约不一致。

\subsection{后续比较与可证伪规则}
\label{sec:future-protocol}

完整主比较需要相同基座、Stage 1 adapter、任务文件、提示、候选数、候选损失设置和
Stage 2 预算。\bone{} 是原始 generation-only 锚点；Random-9 衡量加入候选损失后的
差异；频率匹配随机检验结构画像是否只是关系频率代理；Profile-only 与 Violation-only
分别检验两个信号，ErrorMemory-9 仅作辅助对照。若 SC-CL 不优于频率匹配随机，组合
采样优势的主张必须撤回；还需单独记录输出格式、未列出的答案与 OOV 情况。

计划在 NELL23K 和 FB15K237-2 上使用训练种子 2026--2030，报告逐种子结果与问题级
配对区间；$K\in\{3,9,15\}$、$\rho\in\{0,1/3,2/3\}$ 和
$\tau_s\in\{0.10,0.25,0.50\}$ 用于敏感性检查。WN18RR 与 CoDEx-M 只有完成任务文本、
adapter、bank 和完整评估的对齐审计后才能进入迁移比较。当前 CoDEx-M 已转换官方正例
split 和固定实体对的关系候选，但 Stage 1/Stage 2 推理及方法训练尚未完成；其通过替换
实体构造的官方硬负例不作为关系候选的真值标签。

每个正式运行需记录 bank 的训练来源、question ID 对齐、每行九个互异候选、train-known
过滤结果、持出正例重叠、数据和 adapter 哈希，以及原始输出与候选评分。
未在训练图出现且未被其他证据确认的候选列为“开放世界未判定”，不计入已证实假阴性。
任何 split 泄漏或任务对齐失败均使该运行退出主比较。

\subsection{默认超参数}

\begin{table}[h]
\caption{SC-CL 的默认设置和正式运行前需要登记的字段。}
\label{tab:hyper}
\centering\small
\begin{tabular}{ll}
\toprule
字段 & 默认值 \\
\midrule
候选数 $K$ & 9 \\
均匀混合系数 $\rho$ & $1/3$ \\
结构权重 $\alpha$ & 1.0 \\
B1 违例权重 $\beta$ & 1.0 \\
采样温度 $\tau_s$ & 0.25 \\
候选损失温度 $\tau_c$ & 运行配置登记 \\
候选损失权重 $\lambda_{\mathrm{cand}}$ & 1.0 \\
bank 随机种子 & 2026 \\
训练种子 & 2026, 2027, 2028, 2029, 2030 \\
过滤来源 & train-only \\
\bottomrule
\end{tabular}
\end{table}

\subsection{候选 bank 的最小 schema}

每一行 bank 至少包含：\texttt{question\_id}、\texttt{split}、\texttt{head}、\texttt{tail}、\texttt{gold\_relation}、
\texttt{negative\_relations}、\texttt{negative\_source}、\texttt{filter\_splits}、\texttt{counterfactual\_operation}、
\texttt{counterfactual\_validation}、\texttt{K}、\texttt{uniform\_mix}、\texttt{structure\_weight}、
\texttt{violation\_weight}、\texttt{temperature}、\texttt{seed} 和 sampler 版本。每个候选记录还应包含
\texttt{structure\_similarity}、\texttt{score\_violation}、\texttt{utility}、\texttt{candidate\_score}、
\texttt{gold\_score}、\texttt{score\_gap} 和 \texttt{counterfactual\_valid}。字段名以项目的实际 JSONL
实现为准，论文只要求这些语义信息可审计。

\subsection{审计伪代码}

\small
\begin{verbatim}
for q in train_questions:
    assert q.split == "train"
    P = train_relations - {q.gold} - known_train_relations(q.head, q.tail)
    for r in P:
        s = profile_similarity(phi[q.gold], phi[r])
        v = b1_score_violation(q, r)
        u[r] = alpha * s + beta * v
    negatives = sample_without_replacement(
        P, mixture_softmax(u, rho, tau), K=9)
    write_frozen_bank(q, negatives, manifest=provenance)
\end{verbatim}

该伪代码强调三个不可省略的顺序：先限定训练 split，再过滤实体对已知关系，最后才计算采样分布。正式 wiring 检查还要确认训练 task 的 question ID、正关系文本和候选关系文本保持一致。

\subsection{失败分类}

\begin{table}[h]
\caption{建议的失败分析标签。}
\label{tab:failure}
\centering\scriptsize
\begin{tabular}{p{0.23\linewidth}p{0.65\linewidth}}
\toprule
标签 & 判定规则 \\
\midrule
结构稀疏 & 目标关系或候选关系的画像特征大部分为零，profile 相似度不稳定。 \\
已知关系冲突 & 候选属于当前实体对 train-known 关系，说明过滤或 bank 版本错误。 \\
开放世界不确定 & 候选未在训练图出现，但在外部/测试事实中可能成立，不能当作确定假阴性。 \\
B1 高置信混淆 & 候选 score 高于 gold，且模型对该关系赋予较高置信度。 \\
候选不足 & 过滤后 $|P_q|<9$，必须进入失败清单，不能重复填充。 \\
输出格式异常 & 模型输出不属于题目给出的关系选项。 \\
\bottomrule
\end{tabular}
\end{table}

\subsection{伦理与数据说明}

本文使用本地知识图谱基准和公开模型技术报告作为背景，不新增个人敏感数据。当前正文
只把带有原始预测、配置、manifest 和 hash 的 NELL 单种子结果作为实测；历史 FB
InfoNCE-i0 另行标注，未完成协议不用于性能承诺。正式提交前应补充数据许可、模型许可
证、代码版本和可公开的 bank/预测文件说明。

